\section{Results and Discussion}
\label{sec:results}

The results are reported and discussed per research question, each section closing with its
discussion and the limitations that bound it. Unless stated otherwise, a score is the mean
within-compound Spearman correlation of Equation~\ref{eq:metric}, out of fold over held-out cell lines,
averaged over the \NPanel{} panel compounds and three seeds \cite{spearman1904,dreval}.
% ; each value is
% read from a committed artifact under \code{notebooks/outputs/}, named in \code{results\_numbers.tex}.

\subsection{Prediction from single cells with bulk labels}
\label{sec:results-q1}
\label{sec:rd-q1}

Figure~\ref{fig:panelresponse} shows the label itself, the refitted AUC of each panel compound across
the \NLinesCV{} cross-validation lines, to establish what there is to predict for each compound before
any model is scored. The focus is the spread of each compound's response, since a compound whose lines
all respond alike cannot be ranked. Each row is one compound, the box spans the 25th to 75th
percentile, the whiskers the 5th to 95th, every line is a dot, and the dashed line marks the no-effect
level of 1.0. The rows are ordered by box width. Paclitaxel, gemcitabine, dasatinib and doxorubicin
spread widely below the no-effect level. Erlotinib, sorafenib and etoposide are centred closer to it,
and crizotinib has a narrow box well below it. Imatinib and cisplatin sit close to the no-effect level, and
cisplatin has the narrowest box of the panel.

\begin{figure}[t]
\centering
\includegraphics[width=0.9\textwidth]{panel_response.png}
\caption{\textbf{The label per panel compound.} Refitted AUC (Section~\ref{sec:target}) of each of the
\NPanel{} panel compounds across the \NLinesCV{} cross-validation lines. Box: 25th to 75th percentile,
line: median, whiskers: 5th to 95th percentile, dots: individual lines with vertical jitter only. The
dashed line marks the no-effect level of 1.0, and lower values mean a more sensitive line. Rows are
ordered by box width as a display choice. The panel was not selected on this spread
(Section~\ref{sec:panel}), and the compound with the narrowest spread is part of it.
% Generated by \code{docs/make\_figures.py} (\code{build\_panel\_response}) from the targets file.
}
\label{fig:panelresponse}
\end{figure}

Table~\ref{tab:sweep} reports the score of every configuration, to establish whether any configuration
predicts response and how the per-cell model compares with the ridge control. The focus is the PCA
column, which holds both the highest score and the control. Rows are the six combinations of loss and
density exponent $\alpha$ for the MLP, followed by the ridge regression on cell-line means. Columns are
the two representations, and the highest score is set in bold \cite{hoerl1970,dir}. All six PCA
configurations scored between \SpSweepLo{} and \SpBest{}, the highest score being reached with MAE and
$\alpha = 0.5$, and the ridge regression scored \RidgePca{}. Under the decision rule no configuration
displaced the incumbent, unweighted MAE at \SwPcaMaeZero{}: the \SpBest{} configuration was better on order
but worse on values and spread by more than their seed bands.
% (\code{panel\_metrics.csv})
One configuration (PCA, MSE, $\alpha = 0$) differed between two identical
executions, \SwPcaMseZeroAlt{} against \SpSweepLo{}. The other thirteen entries agreed to four decimals.
% (\code{panel\_execution\_band.csv})
The scoring conventions behaved as expected. Pooled over all (line, compound) pairs rather than within
each compound, the highest-scoring configuration reached a correlation of about \SpPooled{}, and
DrEval's predictor that knows only each compound's average response reached \DrDrugMean{}, so a pooled
score measures mainly the differences in potency between compounds \cite{dreval}. Averaging the
within-compound score per fold, over about 30 lines each, gave a higher score than pooling the held-out
lines in every configuration and seed, so the pooled value reported here is the more conservative of
the two.
% (\code{dreval/bias\_accounting.csv}, \code{panel\_aggregation\_comparison.csv})

\begin{table}[h]
\caption{\textbf{Out-of-fold scores of all configurations.} Mean within-compound Spearman correlation
between predicted and true AUC across held-out cell lines (Equation~\ref{eq:metric}), for the MLP under
each loss (MSE, MAE) and density exponent $\alpha$ (Equation~\ref{eq:density}, where $\alpha = 0$ is
unweighted), and for the ridge regression on cell-line means, on either representation. Five-fold CV
grouped by cell line over \NLinesCV{} lines, \NPanel{} compounds, mean over seeds 42--44. Higher is
better, 0 is no ranking skill. Bold: highest score.
% Read from \code{panel\_per\_drug\_correlation.csv}
% and \code{panel\_ridge\_baseline.csv} (\code{notebooks/4a\_percell\_training.ipynb}).
}
\label{tab:sweep}
\centering
\begin{tabular}{llcc}
\toprule
Model & Loss, $\alpha$ & PCA & scGPT \\
\midrule
MLP & MSE, 0   & \SwPcaMseZero & \SwScgptMseZero \\
MLP & MSE, 0.5 & \SwPcaMseHalf & \SwScgptMseHalf \\
MLP & MSE, 1   & \SwPcaMseOne  & \SwScgptMseOne \\
MLP & MAE, 0   & \SwPcaMaeZero & \SwScgptMaeZero \\
MLP & MAE, 0.5 & \textbf{\SwPcaMaeHalf} & \SwScgptMaeHalf \\
MLP & MAE, 1   & \SwPcaMaeOne  & \SwScgptMaeOne \\
\midrule
Ridge (line means) & --- & \RidgePca & \RidgeScgpt \\
\bottomrule
\end{tabular}
\end{table}

Figure~\ref{fig:labelquality} breaks the highest-scoring configuration down by compound, to test
whether the variation between compounds follows the quality of their labels. The focus is the relation
between the two axes rather than any single compound. Each point is one
compound, placed by the share of its dose--response curves that are flat (logarithmic axis) and by its
score. The compounds at the extremes and the exception are labelled in bold. Scores ranged from
\SpDrugLo{} for cisplatin to \SpDrugHi{} for dasatinib, and the share of flat curves from $0$ to
\PlatinFracNs\,\%. The two were correlated, Spearman $\rho = \FlatRho$ ($p = \FlatP$, $n = 11$):
compounds with more flat curves scored lower, cisplatin having both the largest share and the lowest
score, while crizotinib had no flat curves and a middling score. Across all \NCurves{} curves behind the
panel, the median $R^2$ of the fits was \CurveMedRsq{}, \CurveFracLowRsq\,\% had $R^2 < 0.5$, and
\CurveFracNs\,\% were flat. Cisplatin's median $R^2$ was \PlatinMedRsq{}.

\begin{figure}[t]
\centering
\includegraphics[width=0.85\textwidth]{label_quality.png}
\caption{\textbf{Per-compound score against label quality.} One point per panel compound
($n = \NPanel$). Vertical axis: mean within-compound Spearman correlation, out of fold, of the
highest-scoring configuration (PCA, MAE, $\alpha = 0.5$). Horizontal axis: share of the compound's CurveCurator fits that are flat, i.e.\ cannot reject a
no-effect null ($p > 0.05$), on a logarithmic scale. Bold labels mark the highest (dasatinib) and
lowest (cisplatin) scores and the exception to the trend (crizotinib). The inset reports the Spearman
correlation over the eleven compounds.
% Generated by \code{docs/make\_figures.py} from
% \code{diagnostics/label\_quality\_vs\_performance.csv} and \code{panel\_per\_drug\_correlation.csv}.
}
\label{fig:labelquality}
\end{figure}

The share of the score that reflects general cell-line fragility rather than compound-specific
response was measured next, because LCO cannot remove it \cite{dreval}. Averaged over the six
configurations, the predictions correlated with fragility at \FragRhoPca{} in the PCA arm and
\FragRhoScgpt{} in the scGPT arm, and controlling for it by Equation~\ref{eq:partial} lowered the score
from \FragRawPca{} to \FragPartialPca{} ($-\FragCostPca\,\%$) and from \FragRawScgpt{} to
\FragPartialScgpt{} ($-\FragCostScgpt\,\%$).
% (\code{bias\_accounting.csv})
Figure~\ref{fig:lpo} shows how much a line's identity alone predicts when the protocol allows it to.
The focus is the size of that baseline relative to the scores above. The two bars are DrEval's LPO
baselines that use no expression: the mean response of the line (black) reached a mean within-compound
correlation of \LpoLineMean{}, and the mean response of its lineage (grey) \LpoTissueMean{}.
% (\code{dreval\_lpo\_results.csv})

\begin{figure}[t]
\centering
\includegraphics[width=0.8\textwidth]{lpo_bias.png}
\caption{\textbf{What cell-line identity alone predicts.} Mean within-compound Spearman correlation of
two DrEval baselines under LPO evaluation, in which held-out (line, compound) pairs
belong to lines seen in training: predicting each line's mean response (black, ``which cell line it
is'') and each lineage's mean response (grey, ``lineage only''). Neither uses expression. Five folds of
about 383 held-out pairs over the \NPanel{} compounds. Whiskers are the SD over folds. LPO scores are not
comparable with the LCO scores of Table~\ref{tab:sweep}.
% Generated by \code{docs/make\_figures.py} from
% \code{dreval\_lpo\_results.csv} (\code{scripts/evaluation/dreval\_lpo.py}).
}
\label{fig:lpo}
\end{figure}

Finally, the per-cell model was set against the two line-level controls and the unselected catalogue
\cite{hoerl1970,breiman2001}. The ridge regression on cell-line means scored \RidgePca{} against
\SpBest{}. On DrEval's LCO benchmark the per-compound random forest scored \DrRfNorm{} on the normalised
metric against \DrMlpNorm{} for the MLP, with a fold SD of \DrRfSd{} against \DrMlpSd{}, and both exceeded
DrEval's naive baselines, which score zero on that metric \cite{dreval}. Over all \NDrugsModelled{}
compounds the score was \SpAllDrugs{}, at most \NullGainMax{} above a per-compound constant.
% (\code{panel\_heads\_summary.csv})


\paragraph{Discussion.} Drug response is predictable from these data,
but modestly, unevenly, and not demonstrably from single cells. The score is out of fold over held-out
lines, stable across seeds, and clears DrEval's naive predictors on a protocol this work did not design
\cite{dreval}. Against that, a ridge regression on cell-line means \cite{hoerl1970} reaches the per-cell model, and
a per-compound random forest \cite{breiman2001} matches it externally, so two structurally different line-level models reach the
same score and the single-cell resolution and the network are not what does the work. On the unselected catalogue the signal nearly vanishes, which marks the
panel's score as a property of eleven well-covered compounds rather than of the screen.

\paragraph{What limits the score.} The spread across compounds is explained on the label side: the share
of flat curves predicts the per-compound score without recourse to biology, which agrees with the known
inconsistency of pharmacogenomic labels \cite{haibekains2013,bayer2023curvecurator}. With one label in nine
carrying no dose--response and roughly \NLinesCV{} independent lines \cite{lazic2010,dreval}, the binding constraint is the
supply and quality of labels, not model capacity. The relation is a trend over eleven points,
and crizotinib shows it is not a law. Cisplatin's position at the bottom was anticipated by the
literature on its determinant (Section~\ref{sec:panel}): only one of four ERCC1 isoforms confers repair
capacity and cisplatin resistance, and no antibody distinguishes them \cite{friboulet_ercc1}, so a
gene-level measurement cannot see the functional isoform.

\paragraph{How much of the score is compound-specific.} About a fifth of what either representation ranks is
general cell-line fragility, in nearly the same proportion for both. Under LPO, where a
baseline can use a line's average response, that baseline alone reaches the score of this model. The
LCO protocol used here hides the channel by construction, because a held-out line has no estimated
effect \cite{dreval}. The reported score is therefore real but partly not compound-specific. What this shared component is
was not measured. One candidate, offered as a hypothesis, is proliferation rate: the AUC, like the IC50,
varies with the number of divisions the cells undergo during the assay, and because response is
normalised to an untreated control that keeps dividing, faster-dividing lines appear more sensitive
\cite{hafner}. Testing it would need each line's division rate, which this study did not use.

\paragraph{LCO only.} The fragility channel is invisible under the protocol used, and whether the model
clears a baseline that can use it would require retraining under an LPO split \cite{dreval}.

\paragraph{Why the panel was not selected on response.} Two earlier panels of
this project had been selected on the response values and were discarded: a gate requiring a compound to both
kill and spare cell lines excluded every cytostatic agent by construction, because a compound that
arrests rather than kills never crosses a viability threshold \cite{hafner}; and a literature panel inherited the same
dependency because its candidates had been pre-ranked on that statistic. Selecting on the spread of
one's own labels enriches for compounds that separate these particular lines, flatters every correlation
measured on them, and reads labels of lines that are later held out \cite{dreval}.

\paragraph{About \NLinesCV{} independent labels, not \NCells{} cells.} Every uncertainty in this report is
governed by the number of cell lines (Section~\ref{sec:data}) \cite{lazic2010}, and the work covers
one database, one response score and \NPanel{} compounds. Cross-database integration, toxicity and
leave-tissue-out or leave-drug-out evaluation \cite{dreval} are not attempted.

\paragraph{One label in nine comes from a flat curve, and all labels are weighted alike.} The objective
gives a flat curve the same weight as a well-determined one. Weighting each label by the quality of its
fit, for example by the $p$-value or $F$-statistic CurveCurator reports per curve
\cite{bayer2023curvecurator}, alongside the density weight, would let the model discount such labels.
CurveCurator reports no per-curve variance, so an inverse-variance weight would first need one. This
has not been done \cite{haibekains2013}.

\paragraph{Untuned training hyperparameters.} The learning rate, gradient-clipping threshold, scheduler
patience and factor, early-stopping patience and the batch size of 128 cells are conventional defaults
from the first prototype, neither sourced nor varied. The learning
rate sets where training peaks and the patience when each run ends, and with about a hundred fitting lines
per fold the batch size may be large relative to the number of independent labels. They were left
unchanged rather than chosen after the results were visible, so the comparison is fair, applied
identically to both arms, but not strong. In addition, training on MPS is not bit-reproducible
(Section~\ref{sec:results-q1}), and the check that an MPS embedding agrees with a CPU embedding of the
same cells was never committed.

\subsection{scGPT against PCA}
\label{sec:results-q2}
\label{sec:rd-q2}

The two arms were compared first on the trained models and then on their inputs.
The PCA arm scored higher than the scGPT arm in all six configurations (Table~\ref{tab:sweep}), under
all four scoring conventions --- cell or line level, pooled or per fold
% (\code{panel\_aggregation\_comparison.csv})
--- and after fragility was controlled out. On the inputs
alone, before training, the between-line share of variance, one minus the within-line share defined in
Section~\ref{sec:mil}, was \BtwSharePca{} for PCA and
\BtwShareScgpt{} for scGPT, a between-to-within ratio of \BtwRatioPca{} against \BtwRatioScgpt{}.
Resampling the \NLinesCV{} lines with replacement, both representations on the same resample
(\BtwBootB{} resamples), the difference of the ratios was \BtwRatioDiff{} with a 95\,\% percentile
interval from \BtwRatioDiffLo{} to \BtwRatioDiffHi{}, and PCA was ahead in \BtwBootPcaAheadPct\,\% of
resamples \cite{efron1993}. The lines, not the cells, are resampled because they are the independent
units \cite{lazic2010}.
% (\code{mil/stage0\_input\_ceiling.csv})

Four alternative explanations of the difference were then tested, each on the same folds.
Figure~\ref{fig:labelsupply} shows the first, label supply, to establish whether scGPT's deficit is a
shortage of labels that more labels would close. The focus is the gap between the two arms as the number
of labelled lines grows. Each point is one arm at one budget, the mean over three seeds, with error bars
spanning the lowest and highest seed. Both arms scored higher with more labelled lines. When the labelled lines per fold were thinned from 103 to 25, the
PCA arm lost \CurveGainPca{} and the scGPT arm \CurveGainScgpt{}, and the difference between the arms was
$+\CurveMarginLo$ at 25 and $+\CurveMarginHi$ at 103 lines, the two lying within each other's seed band
at 25 and 50 lines. Replacing the linear head by the
% (\code{panel\_curve\_summary.csv})
$(\CfgHidA,\CfgHidB)$ trunk changed the PCA arm from \ArchLinPca{} to \ArchTrunkPca{} and the scGPT arm
from \ArchLinScgpt{} to \ArchTrunkScgpt{}. The PCA input was about
% (\code{panel\_arch\_summary.csv})
\InputScaleRatio{} times larger in magnitude than the scGPT input.
% (\code{diagnostics/input\_scale.csv})
Rescaling scGPT to the PCA magnitude changed its mean best epoch from \ScaleEpochNative{} to
\ScaleEpochScaled{}, a result whose output was not committed.
% (\code{scripts/evaluation/input\_scale\_test.py})
With fragility controlled out, the highest-scoring
configuration scored \FragBestPca{} in the PCA arm and \FragBestScgpt{} in the scGPT arm.

\begin{figure}[t]
\centering
\includegraphics[width=0.72\textwidth]{label_supply.png}
\caption{\textbf{Score against label supply.} Mean within-compound Spearman correlation of the per-cell
model on PCA (blue) and scGPT (amber), with the labelled lines of each fold's training side thinned to
25, 50 and 75 and at the full 103. Points are the mean over seeds 42--44, error bars the lowest and
highest seed, and the two arms are offset slightly along the horizontal axis so that their bars do not
overlap. All runs used the unweighted MSE configuration, not the highest-scoring one of
Table~\ref{tab:sweep}.
% Generated by \code{docs/make\_figures.py} (\code{build\_label\_supply}) from
% \code{panel/panel\_curve\_summary.csv}; the configuration is read from \code{panel\_curve\_folds.csv}.
}
\label{fig:labelsupply}
\end{figure}

The input-dropout rate was varied last, because it is the one regulariser applied identically to
inputs of different structure \cite{srivastava2014}. Over six rates the PCA arm varied by
\DropRangePca{} and the scGPT arm by \DropRangeScgpt{}. The PCA arm scored \DropBestPca{} at rate 0.02
and \DropShippedPca{} at the rate of \CfgInDrop{} used for all other results.

The size of the gene set made no difference beyond fold-to-fold variation. Over five sizes from
1{,}000 genes to the full transcriptome, the PCA arm's number of compounds beating their per-compound
baseline, out of \NDrugsModelled{}, varied by \SweepRangePca{}, against a fold-to-fold SD of \SweepFoldSd{}, and both representations were flat within error. For scGPT the sweep compared
selection rather than count: at \NHvg{} genes it read about \GenesSeenHvg{} dispersion-selected genes per
cell, in the full-transcriptome variant about \GenesSeenAll{} genes drawn at random under its cap. The
flat result therefore shows only that twice as many random genes do no better than selected ones, not
that scGPT gains nothing from the full transcriptome.
% (\code{notebooks/analysis/qc/verify\_variants.ipynb}, \S9)
% (\code{diagnostics/input\_dropout\_test.csv}, \code{input\_dropout\_sweep\_extra.csv})


\paragraph{Discussion.} scGPT is a viable representation but falls behind PCA here, and the
leading account is a mismatch of granularity rather than a deficit of the embedding. The label is
defined per cell line, so what a model must read is variation between lines. PCA carries about a fifth
more of that structure before any model is trained, which Figure~\ref{fig:umap} illustrates.
scGPT was used frozen, because about \NLinesCV{} labelled lines are too few to fine-tune it and the
project's time did not allow it, and it was
chosen in the expectation that its embedding would carry less cell-line identity, an impression the UMAP
of the atlas gives (Figure~\ref{fig:umap}) and the variance ratio, computed on the model inputs,
confirms. Its authors note that the current pretraining does not
inherently mitigate batch effects \cite{scgpt}, so this reduction is an observation on this atlas, not a
design property of the model. On this label it removes exactly the structure the label varies over. Four rival accounts are eliminated: the margin grows rather than shrinks with more labels,
which contradicts label starvation; capacity lowers both arms; rescaling the input left
the epoch of best validation error almost unchanged, consistent with the Adam update underlying AdamW,
whose step magnitudes are invariant to rescaling of the gradient \cite{kingma2015adam}, although
rescaling the input is not the same as rescaling the gradient; and PCA stays ahead once fragility is removed. Eliminating rivals does not
establish the granularity account, whose positive evidence is one variance ratio, whose difference between the arms survives resampling
the lines, illustrated by one figure. The
margins are conservative with respect to input dropout \cite{srivastava2014}, since the shared rate costs the PCA arm and not
the scGPT arm. This is a statement about evaluation design --- a representation
that carries less cell-line identity on this atlas, scored on a label that is defined per cell line ---
not about scGPT in general.

\paragraph{How much the representations were allowed to see.} The HVG selection, the per-gene standardisation and the PCA rotation were decided together because they
pose one question: how much a fit may see \cite{dreval}. Selecting genes per fold would have produced five gene sets, so comparisons across folds,
representations and the gene-set sweep would differ in genes as well as cells. The lines without
labels were kept in the gene selection, since they carry no label to leak, improve the variance
estimates, and both arms receive the identical set. The per-fold standardisation and rotation were the
one place the arms differ in what is fitted on the data, so their fitting set was restricted to the cells the model
itself is fitted on: the whole training side would have given PCA about \CfgExtraPct\,\% more cells than the
model it feeds, at a cost of about \CfgHoldoutPct\,\% of the cells available to the decomposition. Storing five
fold-specific matrices was rejected, at about a gigabyte and because it would tie the fold design to the
preprocessing. Two asymmetries
remained: PCA decomposes every selected gene while scGPT reads only those in its vocabulary, and PCA is
fitted on this atlas while scGPT is frozen, so the comparison is not like-for-like with respect to
adaptation to the data.

\paragraph{Only one representation is fitted on this atlas.} PCA is fitted on these cells and scGPT is
frozen, so the comparison is not like-for-like with respect to adaptation, and the direct test of the
granularity account, fine-tuning scGPT on the task, has not been run \cite{scgpt}.

\subsection{Heterogeneity learned by a single-cell model}
\label{sec:results-q3}
\label{sec:rd-q3}

The MIL model \cite{dietterich1997,ilse2018} was run over three seeds on both representations and
evaluated by the five stages of Section~\ref{sec:mil}, each on the \QPairs{} (compound, line) pairs of
the panel. The positive control is reported last because it calibrates the others. The first stage,
the input ceiling, is the between-to-within ratio of Section~\ref{sec:rd-q2}, and neither representation
collapsed a line's cells onto one point.
% (\code{notebooks/4b\_mil\_training.ipynb}, \code{mil/q2\_verdict.csv})

The spread stage asked whether the MIL model gives a line's cells more different predictions than the
per-cell model, whose objective penalises exactly that spread (Equation~\ref{eq:mil-decomposition}).
It did so in \QSpreadPca\,\% of pairs for PCA and \QSpreadScgpt\,\% for scGPT, on all three seeds. For
PCA the median within-line SD of the predictions was \QSdMilPcaLo{} to \QSdMilPcaHi{} across seeds,
against \QSdCellPcaLo{} to \QSdCellPcaHi{} for the per-cell model. For scGPT it was \QSdMilScgptLo{} to
\QSdMilScgptHi{}, against \QSdCellScgptLo{} to \QSdCellScgptHi{}.
% (\code{mil/stage1\_spread\_vs\_percell.csv})

The reproducibility stage asked whether that spread is the same in independently initialised models.
Between pairs of seeds, the per-cell predictions within a line correlated with a median $\rho$ of
\QSeedRhoPca{} for PCA and \QSeedRhoScgpt{} for scGPT, and the correlation exceeded the null from
permuting cells within the line in \QSeedSigPca\,\% and \QSeedSigScgpt\,\% of pairs.
% (\code{mil/stage2\_cross\_seed\_agreement.csv})

The confound stage asked how much of that reproducible spread is explained by sequencing depth,
detected genes, mitochondrial fraction and the two cell-cycle scores. The median $R^2$ of these five
covariates was \QConfRawPca{} for PCA and \QConfRawScgpt{} for scGPT, and the adjusted $R^2$, on which the
veto is read, was \QConfPca{} and \QConfScgpt{}. Against the squared cross-seed correlation of
\QReprPca{} and \QReprScgpt{}, that is \QBarPca\,\% and \QBarScgpt\,\% of the veto bar, so the veto did not
fire for either representation.
% (\code{mil/stage6\_confounds.csv})

Figure~\ref{fig:q2instrument} shows the positive control, to establish whether the model can separate
two populations that are known to differ. The focus is the position of the distributions relative to
chance \cite{hanley1982}. Each histogram collects the within-bag AUROC of one arm, PCA in blue and
scGPT in gold, with chance marked by the dashed line. The median was \QAurocPca{} for PCA and
\QAurocScgpt{} for scGPT, and \QAurocSigPca\,\% and \QAurocSigScgpt\,\% of the pairs separated
significantly from the permutation null. The PCA distribution was close to uniform, whereas the scGPT distribution was concentrated at 0
and 1.

\begin{figure}[t]
\centering
\includegraphics[width=0.78\textwidth]{q2_instrument.png}
\caption{\textbf{Positive control of the heterogeneity criterion.} Distribution of the within-bag area
under the receiver operating characteristic curve (AUROC) with which the MIL model separates the cells of
two mixed cell lines whose response is known to differ, computed as the Mann--Whitney statistic divided
by $n_A n_B$ \cite{mannwhitney1947,hanley1982}. One value per (compound, line pair, seed), $17{,}427$ per
representation, is shown as a density histogram for PCA (blue) and scGPT (gold). The dashed line marks
chance (0.5). A value near 1 or 0 means the two populations are ranked apart. Medians: \QAurocPca{}
(PCA), \QAurocScgpt{} (scGPT).
% Read from \code{mil/stage7\_positive\_control.csv}.
}
\label{fig:q2instrument}
\end{figure}

\paragraph{Discussion.} The question remains open. The MIL model
separates a line's cells reproducibly and not as a sequencing artefact, but it does not separate two
lines known to differ: its positive control is barely above chance, so the instrument cannot carry a
negative result and makes the positive one weak. The PCA arm clears the
confound veto at \QBarPca\,\% of the bar, so its within-line signal is not confound-dominated rather than
confound-free, and whether any of the variation is drug-response heterogeneity cannot be decided here. The closest
measurement, the source study's drug screen of sorted subpopulations, found that co-existing
subpopulations differ in drug sensitivity, but it covers two head and neck cancer lines and two
subpopulations, not the cells of every line \cite{scp542}. The criterion can establish whether per-cell predictions
vary within a line reproducibly and not as a sequencing artefact. Whether that variation is
drug-response heterogeneity cannot be established, because no per-cell response was measured.

\paragraph{Design of the MIL model.} A variance term in the loss was rejected
because it would make the spread criterion pass by construction and change objective and architecture at
once. A maximum or quantile aggregator was rejected because it does not force variation either and does
not match an averaged label. This departs from Ilse et al., who, for binary labels under the assumption
that a bag is positive if any of its instances is, call the mean a poor pooling of instance scores
\cite{ilse2018}. Here the label is a continuous average over a line's cells, which the mean matches by
definition. Had the positive control failed, the aggregator was not to be exchanged for one that
passes. Per-line centring of expression, rejected for the per-cell model (Section~\ref{sec:data}), is the
natural representation for this question, and the point to revisit it is a model that pools a line's cells
(Section~\ref{sec:mil}).

\paragraph{The heterogeneity instrument is barely above chance.} A negative answer to the third question
would be uninformative, and little exists to validate a positive one: the only response measurement
below the line level for this atlas is the source study's drug screen of sorted subpopulations of two
head and neck cancer lines \cite{scp542}, and post-treatment single-cell profiles exist for other lines
and compounds \cite{srivatsan2020,tahoe100m} but carry no per-cell response label.
